\chapter{Other works \& achievements}\label{chapter:other}

This chapter opens with a brief overview of other papers I co-authored during my PhD studies. These works either fall outside the thematic scope of this thesis or are part of ongoing research and are therefore omitted from the main text. Nevertheless, as they are relevant to my broader research area, they are included here. The chapter then concludes with a summary of my other achievements, including complementary research activities and outreach initiatives aimed at advancing deep learning.

\section{Publications and preprints}

\begin{itemize}
    \item \textbf{\bibentry{kargin2026sparrta}}
    
    Abstract:
    Visual Foundation Models (VFMs), such as DINO and CLIP, excel in semantic understanding of images but exhibit limited spatial reasoning capabilities, which limits their applicability to embodied systems. As a result, recent work incorporates some 3D tasks (such as depth estimation) into VFM training. However, VFM performance remains inconsistent across other spatial tasks, raising the question of whether these models truly have spatial awareness or overfit to specific 3D objectives. To address this question, we introduce the Spatial Relation Recognition Task (SpaRRTa) benchmark, which evaluates the ability of VFMs to identify relative positions of objects in the image. Unlike traditional 3D objectives that focus on precise metric prediction (e.g., surface normal estimation), SpaRRTa probes a fundamental capability underpinning more advanced forms of human-like spatial understanding. SpaRRTa generates an arbitrary number of photorealistic images with diverse scenes and fully controllable object arrangements, along with freely accessible spatial annotations. Evaluating a range of state-of-the-art VFMs, we reveal significant disparities between their spatial reasoning abilities. Through our analysis, we provide insights into the mechanisms that support or hinder spatial awareness in modern VFMs. We hope that SpaRRTa will serve as a useful tool for guiding the development of future spatially aware visual models.
    \item \textbf{\bibentry{sroka2025ai}}
    
    Abstract:
    Sepsis is a life-threatening condition which requires rapid diagnosis and treatment. Traditional microbiological methods are time-consuming and expensive. In response to these challenges, deep learning algorithms were developed to identify 14 bacteria species and 3 yeast-like fungi from microscopic images of Gram-stained smears of positive blood samples from sepsis patients.
A total of 16,637 Gram-stained microscopic images were used in the study. The analysis used the Cellpose 3 model for segmentation and Attention-based Deep Multiple Instance Learning for classification. Our model achieved an accuracy of 77.15\% for bacteria and 71.39\% for fungi, with ROC AUC of 0.97 and 0.88, respectively. The highest values, reaching up to 96.2\%, were obtained for Cutibacterium acnes, Enterococcus faecium, Stenotrophomonas maltophilia and Nakaseomyces glabratus. Classification difficulties were observed in closely related species, such as Staphylococcus hominis and Staphylococcus haemolyticus, due to morphological similarity, and within Candida albicans due to high morphotic diversity.
The study confirms the potential of our model for microbial classification, but it also indicates the need for further optimisation and expansion of the training data set. In the future, this technology could support microbial diagnosis, reducing diagnostic time and improving the effectiveness of sepsis treatment due to its simplicity and accessibility. Part of the results presented in this publication was covered by a patent application at the European Patent Office EP24461637.1 "A computer implemented method for identifying a microorganism in a blood and a data processing system therefor".
    \item \textbf{\bibentry{pardyl2023complung}}
    
    Abstract:
    Lung cancer is a leading cause of cancer-related deaths, and early diagnosis is crucial for its effective treatment. That is why computer-aided tools have been developed to support particular steps of CT scan analysis, including lung segmentation, suspicious region detection, and patient-level diagnosis. However, none of the previ-ous approaches addressed this process comprehensively. To fill this gap, we introduce CompLung, a comprehensive tool for lung can-cer diagnosis that performs all of the above-listed steps in an end-to-end manner. We have trained the CompLung architecture using the publicly available LIDC-IDRI dataset extended with lung seg-mentation masks obtained from our internal radiologists, which we make publicly available to boost the research on this emerging topic. Finally, we conduct extensive experiments and demonstrate the su-perior performance and interpretability of CompLung compared to existing methods for lung cancer diagnosis.
    \item \textbf{\bibentry{struski2023propml}}
    
    Abstract:
    Partial Multi-label Learning (PML) is a type of weakly supervised learning where each training instance corresponds to a set of candidate labels, among which only some are true. In this paper, we introduce ProPML, a novel probabilistic approach to this problem that extends the binary cross entropy to the PML setup. In contrast to existing methods, it does not require suboptimal disambiguation and, as such, can be applied to any deep architecture. Furthermore, experiments conducted on artificial and real-world datasets indicate that ProPML outperforms existing approaches, especially for high noise in a candidate set.
    \item \textbf{\bibentry{pardyl2022automating}}

    Abstract:
    As the leading cause of cancer-related mortality, lung cancer is responsible for more deaths than colon, breast, and prostate cancer put together. Screening with low-dose computed tomography detects cancer at an early stage and reduces mortality. However, it requires the tedious work of radiologists to obtain malignancy scores, which additionally are very subjective. That is why many researchers worked on methods automating lung cancer classification, usually using the publicly available LIDC-IDRI dataset for training. However, most of those methods consider only node-level classification and provide poor results for patient-level diagnosis. In this paper, we fill this gap by introducing an end-to-end methods with a CT scan on the input and the patient-level diagnosis on the output. We consider three approaches for three different data regimes to examine how stronger and weaker supervision influences the model performance.
\end{itemize}

\section{Patents}

\begin{itemize}
    \item \textbf{European Patent application EP24461637.1} ``A computer implemented method for identifying a microorganism in a blood and a data processing system for therefore''. European Patent Office, 2024, patent pending.
\end{itemize}

\section{Ongoing research}

\begin{itemize}
    \item \textbf{Reinforcement learning for visual exploration in real-world environments.} A. Pardyl, J. Mikołajczyk, J. Matusiak, H. Zientata, D. Gaweł, I. Alwasiak, B. Zieliński, M. Cygan, M. Wołczyk.
    
    Abstract: This research project is a direct follow-up of our FlySearch~\cite{pardyl2025flysearch} benchmark.
    We aim to develop a reinforcement learning framework for training vision-language model-based agents to navigate and explore real-world environments.
    Building upon the GRPO algorithm~\cite{shao2024deepseekmath}, we will fine-tune 4-11B parameter models on a simplified FlySearch benchmark (FlySearch-RealTime), and a collection of Gaussian splat scenes generated from real-world flight data. Such a setup will allow us to
    train the model on diverse navigation tasks (procedurally generated with FlySearch-RealTime), while minimising the sim-to-real gap by including real-world visual data in the training set. We will then evaluate the performance of the trained models on the original FlySearch benchmark, and in a real-world setting with a physical UAV.

    \item \textbf{Self-supervised concept discovery.} A. Pardyl, S. Gairola, S. Rao, B. Schiele, B. Zieliński, D. Rymarczyk.
    
    Abstract: The idea of this project is to create a model that will learn a set of discrete, spatially grounded concepts using only self-supervised learning. Each of the learned concepts represents a single well-defined visual feature, e.g. a semantic part of an object or a texture pattern. Inspired by PDiscoFormer~\cite{aniraj2024pdiscoformer}, the project builds upon the DINOv2 family~\cite{oquab2024dinov2}, extending it with a dedicated attention-based concept learning pathway and a VAE-based quantization stack. The resulting embedding is a drop-in replacement for the ViT CLS token, providing an interpretable backbone for multiple downstream tasks.
    
\end{itemize}

\section{International research internships}
\begin{itemize}
    \item \textbf{Max Planck Institute for Informatics}, June--September 2026, Computer Vision Group, supervised by Prof. Dr. Bernt Schiele.
\end{itemize}

\section{Scientific grants}
\subsection{As a principal investigator}
\begin{itemize}
    \item \textbf{National Science Center Preludium} ``Where to look next - guiding active visual exploration with
    internal model uncertainty'', 2024-2025. Grant number: 2023/49/N/ST6/02465.
    \item \textbf{Foundation for Polish Science START scholarship}, 2026-2027, Grant number: 59.2026.
\end{itemize}

\subsection{As a participant / co-investigator}
\begin{itemize}
    \item \textbf{National Science Center Sonata bis} ``Sustainable computer vision for autonomous machines'', 2024-2027. Grant number: 2023/50/E/ST6/00469. Principal investigator: Bartosz Zieliński.
    \item \textbf{Foundation for Polish Science Proof of Concept} ``Intelligent Fight Against Sepsis - Innovative Use of Artificial Intelligence in Rapid Microbiological Diagnosis of Sepsis'', 2024-2025. Grant number: FENG.02.07-IP.05-0068/23. Principal investigator: Monika Brzychczy-Włoch.
    \item \textbf{European Commission Horizon Europe} ``European Lighthouse of AI for Sustainability (ELIAS)'', 2024-2027. Grant number: 101120237. Principal investigator: Massimiliano Pontil.
    \item \textbf{Foundation for Polish Science Team-Net} ``Bio-inspired artificial neural networks'', 2022-2023. Grant number: POIR.04.04.00-00-14DE/18-00. Principal investigator: Jacek Tabor.
    \item \textbf{National Center for Research and Development Fast Track} ``X-rAI: Diagnostic viewer for computer-assisted radiology using Artificial Intelligence'', 2021-2023. Grant number: POIR.01.01.01-00-1666/20. Principal investigator: Zbisław Tabor.
\end{itemize}

\section{Acting as a reviewer}
\begin{itemize}
    \item {Annual Conference on Neural Information Processing Systems (NeurIPS) 2026}
    \item {CVPR 2025 Workshop on Computer Vision for Drug Discovery (CVDD) 2025}
    \item {IEEE/CVF Winter Conference on Applications of Computer Vision (WACV) 2025}
    \item {ACM International Conference on Knowledge Discovery and Data Mining (KDD) 2024}
    \item {IEEE/CVF Winter Conference on Applications of Computer Vision (WACV) 2024}
\end{itemize}

\section{Doctoral schools and scientific events}
\subsection{As an organiser}
\begin{itemize}
    \item Machine Learning Summer School on Reliability \& Safety (MLSS$^{R\&S}$ 2026), Kraków, Poland
    \item ELLIS Doctoral Symposium (EDS) on Robust AI 2025, Warsaw, Poland
    \item Machine Learning Summer School on Drug and Materials Discovery (MLSS$^{D}$ 2025), Kraków, Poland
    \item Machine Learning Summer School on Applications in Science (MLSS$^{S}$ 2023), Kraków, Poland
\end{itemize}

\subsection{As a presenter}
\begin{itemize}
    \item ML in PL Conference 2025, Warsaw, Poland, ``FlySearch: Exploring how vision-language models explore''
    \item ML in PL Conference 2024, Warsaw, Poland, ``AdaGlimpse: Active Visual Exploration with Arbitrary Glimpse Position and Scale''
    \item SFI IT Academic Festival 2024, Kraków, Poland, ``AI for autonomous agents''
\end{itemize}

\subsection{As a participant}
\begin{itemize}
    \item Mediterranean Machine Learning (M2L) Summer School 2025, Split, Croatia
    \item International Computer Vision Summer School (ICVSS) 2025, Scicli, Italy
    \item ELIAS-ELLIS-VISMAC Winter School 2025, Brunico, Italy
    \item ELLIS Doctoral Symposium (EDS) on AI \& Sustainability 2024, Paris, France
    \item Mediterranean Machine Learning (M2L) Summer School 2023, Thessaloniki, Greece
\end{itemize}

\section{Other activities}
\begin{itemize}
    \item Member of the Audit Committee of the PhD Student Association of the Jagiellonian University (2025-2026)
    \item Member of the PhD Student Council of the Doctoral School of Exact and Natural Sciences of the Jagiellonian University (2024-2026)
    \item PhD Student Representative to the Standing Rector's Committee on Statutory and Legal Affairs (2022-2026)
    \item PhD Student Representative to the Council of the Faculty of Mathematics and Computer Science of the Jagiellonian University (2025-2026)
\end{itemize}